\section{Project Constraints}
\textbf{High-level constraints.} Limited compute (a single Windows machine) restricted quantum-circuit simulation to small qubit counts; the public dataset has very few fraud cases (0.17\% of transactions), which limits statistical power and motivates paired multi-seed designs; there was no physical QKD hardware; time was limited to the project period (repository history spans 2026-08-25 to 2026-10-07 in its recorded commits).

\textbf{Implementation constraints.}
\begin{itemize}[noitemsep]
\item Native \texttt{liboqs} could not be built (no CMake or MSVC toolchain, which was deliberately not installed); the compiled \texttt{pqcrypto} package was used instead, and no liboqs measurements exist.
\item The reference SLH-DSA implementation is computationally expensive (seconds per signature), so its repetition counts are small (3 for key generation and signing).
\item Quantum simulation overhead limited the Qiskit BB84 trials to 500 bits per trial (protocol simulator: up to 2000 bits).
\item IBM Quantum open-plan access limits runtime and queue time, so the hardware campaign used small circuits and five jobs within one session.
\item Simulation assumptions: both BB84 implementations are software simulations with different noise models; no real fraud-bank data were available.
\item The \texttt{pqcrypto} build is compiled but not necessarily the most optimized code path; absolute latencies are machine-specific.
\end{itemize}
